\section{Lossless Coding}
Lossless coding maps source data to a reversible bitstream: decoding reproduces every original symbol exactly. Compression comes only from statistical structure and repeated patterns, not from discarding information.
\subsection{Definitions}
\begin{itemize}
    \item \textbf{Alphabet}: $\mathcal{X} = \left\{x_1, \cdots, x_M\right\}$ Symbols to encode
    \item \textbf{Code}: Application $\mathcal{C} : \mathcal{X} \to \left\{0,1\right\}^\star$ set of finite-length bit strings, can be Fixed or Variable length.
\end{itemize}
\subsection{Code Types}
\subsubsection{Fixed-Length Code}
A fixed-length code assigns the same number of bits to every source symbol. It offers simple parsing and constant per-symbol rate, but cannot assign shorter descriptions to frequent symbols.
Good: straight forward parsing of the codewords.\\
Bad: Very large rate $R = \log_2 L$, we are assuming all symbols are equally probable.\\
Needed bits: $\left\lceil \log_2 M \right\rceil$
\subsubsection{Variable-Length Code}
A variable-length code assigns different codeword lengths to different symbols. Its rate is the expected codeword length $\bar{L}$, measured in bits per source symbol, so frequent symbols should receive shorter codewords.
\begin{itemize}
    \item Statistical Approach: we make use of source entropy $H(x)$, to build an optimal Huffman code. We could use also suboptimal aritmethic codes, which scale very well.
    \item Euristic Approach: if we don't know data statistics, universal coding (LZW)
\end{itemize}
We define $l_i$ as the length of the codeword $c_i$.\\
We need also two conditions:
\begin{itemize}
    \item Decodability condition / prefix condition (each word is not a prefix for other words)
    \item Non-equiprobable symbols
\end{itemize}
How to write this in math terms?
$$f : x_i \to c_i \qquad \qquad \text{avg. length }\bar{L} = \sum p_i l_i$$
We want to minimize the average length.
\subsection{Variable-length codes theorems}
\subsubsection{McMillan's Theorem}
Decodable codes do not improve performance w.r.t. instantaneous/prefix codes.\\
The best possible prefix code = the best possible decodable code. We have to focus on prefix codes.
\subsubsection{Kraft's Inequality}
$$\sr{}{i}2^{-l_i} \leq 1 \qquad \Longleftrightarrow \qquad \exists \ \text{instantaneous code with lengths } \{l_1, \cdots, l_M\}$$
If the equality is verified, we say the code is \textbf{complete}, otherwise we can try to push more from the code.\\
We can associate a binary code to a binary tree.
\subsubsection{Proof for Kraft's $\Rightarrow$ (necessity)}
Build a binary tree with $L_{\max} = \max(l_i)$ depth.\\
We remove nodes which have prefix equal to a codeword, then each removed subset is disjoint from the others:
$$\text{Removed nodes }: \sr{}{i}2^{L_{\max} - l_i} \leq 2^{L_{\max}}$$
\subsubsection{Proof for Kraft's $\Leftarrow$ (sufficiency)}
Greedy Construction Algorithm:
\begin{itemize}
    \item Sort Length $l_1 \leq \cdots \leq l_N$
    \item $\forall$ step $k$: pick available node at depth $l_k$
    \item Mark descendants as forbidden
\end{itemize}
Finally: there is no way to improve it.
\subsection{Recall: Information Theory}
\subsubsection{Information}
Self-information measures how surprising one event is. It is measured in bits when the logarithm has base two: rare events carry more information than probable events.
$$I(x_i) = - \log_2 p(x_i) = -\log_2 p_i$$
Joint Information:
$$I(x_i,x_j) \leq I(x_i) + I(x_j)$$
Equality only if symbols are independent.
\subsubsection{Source Entropy}
Source entropy is the average self-information generated per symbol by a memoryless source. It is measured in bits per symbol and is the fundamental lower bound approached by lossless codes matched to the source statistics.
$$H(X) = \sr{}{i}p_iI(x_i) = -\sr{}{i}p_i\log_2 p_i$$
\subsubsection{Max Entropy Distribution}
We want $p^{\star}$ such that:
$$p^{\star} = \arg\max_{p}\sr{M}{i=1}p_iH(x_i) \qquad \text{with }\sr{M}{i=1}p_i = 1$$
Using Lagrange Multipliers:
\begin{itemize}
    \item we consider a function $f:\mathbb{R}^n \to \mathbb{R}$
    \item then we build a function $J(x, \lambda) = f(x) + \lambda\cdot \varphi(x)$ 
    \item then we set $\forall i$ the $\nabla J(x_i, \lambda) = 0$
\end{itemize}
In this case:
\begin{itemize}
    \item $f(p) = -\sr{M}{i=1}p_i\log_2(p_i)$
    \item $\varphi(p) = \sr{M}{i=1}p_i - 1$
\end{itemize}
Thus the gradient is:
$$\frac{\partial J}{\partial x_i} = -\left({\log_2 e} + \log_2 p_i\right)+\lambda \qquad \qquad \qquad \frac{\partial J}{\partial p_i}(p^\star) = \varphi(p) = 0$$
By imposing =0 to the $x_i$ component of gradient:
$$\log_2p_i^\star = \lambda - \log_2 e$$
\subsubsection{Joint Entropy}
Joint entropy is the average information needed to describe two random variables together. It includes both their individual uncertainty and any statistical dependence between them.
$$H(X,Y) = -\sr{}{i,j}p_{ij}\log p_{ij}$$
\subsubsection{Conditional Entropy}
Conditional entropy $H(X|Y)$ is the uncertainty that remains about $X$ after $Y$ is known. In coding terms, it is the ideal average rate for coding $X$ when the decoder also has $Y$ as side information.
$$H(X|Y) = \sr{}{j}p_jH(X|Y=Y_j)$$
Using chain rule:
$$H(X,Y) = H(Y) + H(X|Y) = H(X) + H(Y|X)$$
